\section{The construction at fixed lattice spacing}
\label{sec:wilson}

This is the strongest constructive statement in the program. It is also
the one most often misread, so it is stated here with its hypotheses,
its interval, and an explicit list of what it does not assert.

Throughout this section $\varepsilon$ denotes the temporal mesh. It is
\emph{not} the planar variable $\tau$ of \eqref{eq:tau}; the two are
unrelated and the symbols are kept distinct deliberately.

\subsection{The theorem}

\begin{theorem}[Infinite-volume Wilson band and literal source frame]
\label{thm:wilsonband}
Under the hypotheses listed below, the Perron-normalized Wilson transfer
operator has a common strong limit along open-box and centered-periodic
exhaustions, in the transported product representation. The limit is a
positive contraction fixing the vacuum; it differs from the free product
in norm by at most $1/998$, and its norm on the non-vacuum subspace is
at most $\tfrac45 + \tfrac1{998}$.

On the coupling interval
\begin{equation}
  \abs{u} \;\le\; \frac{u_*}{10\,022\,400\,000\,N},
  \label{eq:wilson-interval}
\end{equation}
the complete infinite-volume physical odd band is isomorphic to
$\ell^2(\Z^3) \otimes \C^3$. The projected literal plaquette-source
synthesis is \emph{onto} that entire band, with Gram bounds
\begin{equation}
  \tfrac{9}{16}\,I \;<\; \mathcal{G} \;\le\; \tfrac{81}{64}\,I .
\end{equation}
The Euclidean Osterwalder--Schrader band and the quantum GNS band
coincide, intertwining the physical transfer and the sources.
\end{theorem}
\gid{RESULT:WILSON\_INFINITE\_PHYSICAL\_BAND}

\noindent Recorded \texttt{proven} with \texttt{analytic} evidence;
\Tthree{} for whole-statement machine certification, with its component
estimates checked across several suites.

\subsection{Hypotheses}

\begin{enumerate}
  \item All calibrated kinetic, positive Wilson-kernel, incidence and
    common-disk hypotheses of the weighted-activity and cardinality-chart
    theorems.
  \item Fixed $\SU(N)$ with $N \ge 3$, and fixed \emph{spatial} lattice
    spacing. Compatible induced open boxes, or centered periodic
    exhaustions, carrying the same local data.
  \item Temporal mesh
    $0 < \varepsilon \le \varepsilon_0 \le \log(5/4)/(4\gamma)$, with
    $m = \lceil \log(5/4)/(\gamma\varepsilon) \rceil$, and the full
    neutral physical odd free-window premise.
  \item $\gamma = C_F/2 = (N^2-1)/(4N)$, the \emph{exact} calibrated
    fundamental-link energy. This is a hypothesis about an identity, not
    a gap lower bound one is free to weaken: substituting a smaller
    number invalidates the conclusion rather than weakening it.
  \item Literal odd sources
    $\mathcal{O}_p = (\chi_p - \bar\chi_p)/\sqrt2$, with free
    orthonormal plaquette vectors of norm at most $\sqrt2\,N$.
  \item The calibrated representation-uniform kinetic-window mesh
    premise, whose $N$-dependent admissible temporal threshold is the
    previously established one. Every estimate is uniform in volume and
    in meshes satisfying these premises, and the common coupling
    interval depends on fixed $N$ and is nonzero.
\end{enumerate}

\noindent The interval \eqref{eq:wilson-interval} unpacks as
$u_0 = u_*/1252800000$ and $u_1 = u_0/(8N)$, with the accompanying
constants $\eta_{\mathrm{act}} = 1/2500$, $\epsilon = 1/998$,
$g_* = 1024/15625$ and $q_0 = \tfrac45 + \tfrac1{998} < 1$. The bound
$\eta_{\mathrm{act}} \le 1/2500$ on the integrated generator norm at
weight $1 + \log 2$ is a \emph{hypothesis} of the source
specialization, not a conclusion of the theorem; it is listed here to
prevent the reading in which the theorem is taken to establish it.

\subsection{Why totality is the substance}

The clause that carries the physical content of
Theorem~\ref{thm:wilsonband} is that the source synthesis is \emph{onto}
the band, not merely bounded into it.

Break~\ref{sec:break4} is the reason. A local probe family can decay at
one rate while the true spectrum contains a much slower state the probe
cannot see; the three-dimensional model of Proposition~\ref{prop:break4}
exhibits a factor of twenty-one between the observed rate and the actual
gap. A construction that produced a band and a decay estimate, without
totality, would prove nothing about the spectrum.

Totality here is a conclusion. The tagged source norm controls the
entire synthesis operator; close projections and a Neumann inverse
establish surjectivity in infinite dimension; and reflection-positive
history completion contains the whole band by selective source
cyclicity. Fine time steps are the positive $m$-th root of the block.

\subsection{What Theorem~\ref{thm:wilsonband} does not assert}

\notasserted{Global equality of the Osterwalder--Schrader and full
quantum high-energy spaces. Global transfer injectivity. Any spatial
continuum limit. Extension beyond $\SU(N)$, $N \ge 3$. The positive
energy gap is in electric-time units, and its conversion to physical
units is not part of the statement.}

Most importantly, and separately from the list above: the interval
\eqref{eq:wilson-interval} is a neighbourhood of $u = 0$, and the
continuum limit is taken along a trajectory on which $u \to \infty$.
Theorem~\ref{thm:wilsonband} therefore supplies no control whatever on
the continuum limit. This is Break~\ref{sec:break5}, it is structural,
and no improvement of the constant $10\,022\,400\,000$ addresses it.
A version of exactly that inference was made in this program's own
drafts and retracted; see \S\ref{sec:register}.

\subsection{The relation to the Hamiltonian carrier}

Theorem~\ref{thm:wilsonband} is a statement about the Wilson transfer
operator. The Hamiltonian carrier of Part~\ref{part:exact} is a
different object, spectrally identified with the band only on momentum
patches away from zero.

The obstruction at zero momentum is proved, not merely encountered, and
it is a statement about \emph{chains}:
\begin{enumerate}
  \item every vector in the image of the cube boundary map has exactly
    zero zone-average component;
  \item the zero-momentum carrier content is therefore carried entirely
    by the three wrapping sheets of \S\ref{sec:carrier};
  \item a sheet's normalized overlap with any fixed local observable is
    exactly $1/L$;
  \item compact two-cycles on open windows are exactly cube boundaries
    --- checked directly for small windows, and Alexander duality in
    general.
\end{enumerate}
The repository's own prose summarizes (i)--(iv) as ``every finitely
supported exactly-flat operator is $\Gamma$-blind''. That phrasing is a
convenient shorthand and this document does not adopt it: the four
statements are about chain vectors, and read literally as a claim about
arbitrary operators the shorthand is false. What follows from (i)--(iv)
is the thing actually needed --- that an interpolator for the
zero-momentum protected level must leave the exactly-flat local algebra,
coupling either through in-sector directions that are flat-degenerate at
$\Gamma$ anyway, or through out-of-sector multi-plaquette content.

\subsection{A prediction, sharper than the measurement it explains}

The published few-percent overlap of a local operator with the lightest
state is not a constant: it degrades toward the continuum. For a local
lattice source whose lowest-dimension continuum operator has dimension
$d$, the overlap law is
\begin{equation}
  Z_0(a) \;\sim\; \frac{\abs{f_d}^2\, a^{2d-3}}{2M},
  \label{eq:overlaplaw}
\end{equation}
which at $d = 4$ gives $a^5$ and so reproduces the published
measurement --- the agreement that makes \eqref{eq:overlaplaw} worth
recording at all.

But $a^5$ is the wrong exponent for \emph{this} carrier. The C-odd
carrier's leading continuum operator is $d^{abc}F^aF^bF^c$, of dimension
six, so the applicable exponent is $a^{2\cdot 6-3} = a^9$. The
extrapolation is brutal: $Z_0 \approx 1.2\times10^{-6}$ at
$a\sqrt\sigma = 0.10$ and $2.4\times10^{-9}$ at $0.05$. Still a power,
so formally survivable, and a serious practical obstacle to measuring
this state with a local operator.

The same law is the argument for smearing, quantitatively: at fixed
physical radius $R$, $Z_0 = \abs{f_R}^2 R^{2d-3}/(2M)$ is
$a$-independent. A smeared basis is structural here, not a convenience.

\begin{scopenote}
Equation~\eqref{eq:overlaplaw} is asserted, not proved here. What is
checked is the dimension arithmetic and the agreement of its $d = 4$
case with the published $a^5$. An earlier version of this program's
notes quoted $a^5$ for the carrier itself; that is corrected above.
\end{scopenote}

\begin{scopenote}
The out-of-sector dressing, the transfer-matrix identification at zero
momentum, and everything non-perturbative remain open. A particle
statement lives at fixed nonzero momentum, where no collapse occurs; the
$L^{-2}$ statement in the literature is a minimal-grid-momentum
statement governing the approach to $\Gamma$, not a fixed-momentum one.
\end{scopenote}
